% ============================================================================
% Supplementary Material for:
%   "Composable, Knowledge-Driven Propagation for Complex Interdependent
%    Services with a Flow-Allocation Module Validated Against Hydraulic
%    Simulation" (Curaba & Grimaz, Complex Networks & Their Applications)
%
% Content moved here from the main paper's appendices (2026-07-17) to meet
% the venue's 12-page cap. Distributed alongside the released code, not
% submitted as part of the proceedings paper. Sections are numbered S1-S8
% and referenced from the main text as "Supp. Mat. SN".
%
% ORDER (reset 2026-08-17): the sections now follow the main paper's own
% narrative — setup (S1 importer, S2 protocol, S3 inputs), then results
% (S4 extended validation, S5 ablations, S6 sensitivity), then what did not
% work (S7), then cost (S8). Any renumbering here MUST be mirrored in every
% "Supp. Mat. SN" string in complenet.tex — those are plain text, not \ref.
%
% STYLE: setup + numbers + take-home message. Tables carry the evidence;
% prose states the design and the conclusion, and stops.
%
% BUILD: pdflatex FlowAllocationModuleForHydraulic-supplementary && bibtex FlowAllocationModuleForHydraulic-supplementary
%        && pdflatex FlowAllocationModuleForHydraulic-supplementary (x2)
%
% Same anonymization hard constraint as the main paper: "Aqueduct A/B/C"
% only, never real utility names or locations.
% ============================================================================
\documentclass[11pt]{article}

\usepackage[a4paper,margin=2.8cm]{geometry}
\usepackage{amsmath,amssymb,amsfonts}
\usepackage{booktabs}
\usepackage{graphicx}
\usepackage{xcolor}
\usepackage{url}
% hyperref LAST, same reasoning as the main paper: `url` typesets an address but
% creates no link. Here it mainly makes the many Sec./Table cross-references
% clickable. `hidelinks` leaves print output unchanged.
\usepackage[hidelinks]{hyperref}
\urlstyle{same}

\newcommand{\worstof}{\mathrm{worst\_of}}
\newcommand{\bestof}{\mathrm{best\_of}}
% Supplementary sections are S1, S2, ...
\renewcommand{\thesection}{S\arabic{section}}
\renewcommand{\thetable}{S\arabic{table}}
\renewcommand{\thefigure}{S\arabic{figure}}

\title{Supplementary Material\\[2mm]
\large Composable, Knowledge-Driven Disservice Propagation for Complex Interdependent
Services with a Flow-Allocation Module Validated Against Hydraulic Simulation}
\author{Cristian Curaba \and Stefano Grimaz\\
\normalsize University of Udine, Udine, Italy}
\date{}

\begin{document}
\maketitle

\noindent This document specifies the setup behind the main paper's results
(Sec.~\ref{app:importer}--\ref{app:inputs}), reports them in full
(Sec.~\ref{app:extended}--\ref{app:sensitivity}), records what did not work
(Sec.~\ref{app:negative}), and measures cost (Sec.~\ref{app:cost}). Main-paper
sections are cited as ``main paper, Sec.~$n$''. The critical-class detection
figures share the main paper's Table~1 with the per-family FMS and are not
repeated here.

% ============================================================================
\section{The EPANET Importer in Full}\label{app:importer}

The main paper's Sec.~4 summarizes the import decisions that carry fidelity.
This section specifies every stage. The parameterization actually used is in
Sec.~\ref{app:repro}.

\emph{Structural mapping.} Reservoirs and tanks both become source nodes, with
supply set by \emph{kind}. A reservoir is a fixed-head body whose outlet pipes
limit delivery, so it becomes an \emph{unbounded} source --- matching how the
WNTR solver treats it, and keeping multi-source failure comparable against the
oracle. A tank's supply is bounded instead: it becomes a \emph{rate-limited} source
whose ceiling is its \emph{nominal delivered outflow} --- the flow observed on its
incident links in one solve at nominal demand. That bound is a \emph{rate}: nothing about it runs out.
A tank's stored volume enters the model as a separate, temporal attribute, its usable volume divided by downstream demand becoming
a \emph{backup duration} in hours, so the reserve is imported as a countdown
that starts when the tank's own feed is cut rather than as extra flow. Node
functionality then scales the rate: a source at level $2$ of $3$ injects a
proportional fraction of its ceiling. Demanding junctions are served
nodes; pumps and valves become inline \textsc{Requisite}-typed nodes, so their
operating state gates downstream flow through the same logical propagation with no
special-cased device logic; pipes become capacitated, directed edges. Junction
demand is taken at the single pattern hour of maximum total consumption.

The tank rule has one exception, and it is load-bearing. A tank reachable from a
reservoir only through a pump is not a store the network draws down but a
\emph{pass-through} the pumps keep full, and its throughput is set by its pipes,
not its geometry. We detect this structurally --- a tank outside the component
every reservoir reaches through pipes and valves alone --- and give such tanks
incident-pipe capacity instead. Without it C-Town, whose seven tanks are all
pump-fed, reads as mass starvation: mean FMS $0.03$ against $0.91$ with the rule.
Which of the two tank behaviours of Sec.~\ref{app:negative} applies to a network
follows from the same distinction.

\emph{Orientation from simulation.} A \texttt{.inp} file's declared pipe
direction is one arbitrary operating snapshot; the true direction can differ and
even reverse under another operating point. We therefore run many steady-state solves
and read each pipe's direction from the \emph{sign} of its flow.

Three cases arise, separated by two velocity thresholds. A link whose fastest
reading in either direction never clears a noise floor ($10^{-4}$\,m/s) carries
no signal: it takes a topology-only (breadth-first) orientation guess and the
$1$\,m/s fallback velocity. A link clearing the floor in \emph{both} directions
is genuinely two-way and becomes two opposed directed edges, each carrying the
pipe's whole capacity. The third case is the one worth stating: a link
clearing the floor in only \emph{one} direction, but at under $0.3$\,m/s ---
below the self-cleansing velocity a real main is designed for, where confidently
one-way trunk flow runs at $0.9$--$4$\,m/s --- is \emph{also} split. The
asymmetry is deliberate: near a near-zero equilibrium the solver's sign is as
likely noise as signal, and forcing the wrong one can silently disconnect every
junction downstream, whereas splitting a pipe that really was one-way costs an
unused reverse capacity. We tried the stricter alternative --- an absolute-flow
threshold plus a dominance ratio to pick a winner --- and reverted it: it fixed
one noisy leaf pipe on the largest aqueduct but scored materially worse across
the benchmark, severing genuine near-balance points.

\emph{The solves.} \emph{(1)~A demand solve} raises every junction's demand
together to several times nominal --- as high as the network still converges,
stepping down adaptively. \emph{(2)~Contingency solves} close one link at a
time: the top tenth of pipes by diameter and every pump, plus a bounded random
sample of the rest, re-solving at nominal demand and undoing each closure before
the next. ``Closing'' sets EPANET status to \textsc{closed}, the same operation
the validation uses for a break. The second family is what reveals
\emph{reversal}: a backup pipe sits nearly idle in normal operation and shows its
working direction only once an upstream link fails and flow reroutes onto it. On
Net3, a $203$\,mm backup main carries almost no flow under demand escalation
alone, but a clear reverse flow once an upstream pipe is closed.

\emph{Capacity from a uniform design velocity.} Diameter alone does not fix
usable capacity, but a diameter and an assumed velocity do: capacity is
$\tfrac{\pi}{4}d^2 v_{\text{design}}$ with $v_{\text{design}}=2.5$\,m/s, the
standard water-main design velocity (an exposed import parameter). We initially
sized capacity per-pipe from the \emph{peak} velocity each pipe reaches across
the solves above; the ablation of Sec.~\ref{app:ablation} found this does no
better, so we adopt the simpler, reproducible constant.

\emph{Skeletonization to a node budget.} Real utility exports exceed practical
importer budgets by an order of magnitude, so the pipeline can reduce to a target
node count via WNTR's demand-conserving skeletonization --- branch trimming plus
series and parallel pipe merges below a diameter threshold, searched for the
smallest threshold meeting the budget. It is bypassed throughout the main paper's
Sec.~5, which imports every network in full.

% ============================================================================
\section{Experimental Setup and Protocol}\label{app:repro}

All numbers in the main paper come from one harness
(\texttt{scripts/validate\_faithfulness.py} in the released code) on a single
machine. This section records every parameter it uses.

\emph{Software.} Python 3.12.3; WNTR 1.5.0 (driving EPANET~2.2); NetworkX 3.6.1;
NumPy 2.5.1. All randomness flows through seeded \texttt{random.Random}
instances, and container iteration is explicitly sorted wherever order can
influence a solver, so runs reproduce across processes and machines.

\emph{Networks.} Three public EPANET examples, three anonymized Italian aqueduct
exports, and two public medium benchmarks; all imported in full, no
skeletonization.

\begin{center}
\begin{tabular}{lrrrrrr}
\toprule
 & junctions & demand $>0$ & pipes & pumps & tanks & reservoirs\\
\midrule
Net1 & 9 & 8 & 12 & 1 & 1 & 1\\
Net2 & 35 & 32 & 40 & 0 & 1 & 0\\
Net3 & 92 & 59 & 117 & 2 & 3 & 2\\
Aqueduct A & 417 & 410 & 437 & 0 & 0 & 2\\
Aqueduct B & 439 & 424 & 499 & 1 & 3 & 1\\
Aqueduct C & 607 & 601 & 674 & 0 & 1 & 1\\
Modena & 268 & 245 & 317 & 0 & 0 & 4\\
C-Town & 388 & 334 & 429 & 11 & 7 & 1\\
\bottomrule
\end{tabular}
\end{center}

\emph{Import parameterization} (Sec.~\ref{app:importer}). Demand snapshot:
\texttt{peak\_hour}, fixed onto every junction before any solve. Capacity
(adopted default): $\pi/4\,d^2\,v_{\text{design}}$, $v_{\text{design}}=2.5$\,m/s,
no solve. The retained per-pipe drill (Table~\ref{tab:capmethod}) instead takes
the maximum peak velocity over (a)~a demand solve at the highest multiplier
$\le 8\times$ the network still converges at and (b)~contingency closures, giving
$\pi/4\,d^2 \min(v_{\mathrm{peak}}\!\times\!2,\ 3\,\text{m/s})$. Both methods
still run the solves for \emph{orientation}: contingency discovery closes the
top-tenth-by-diameter pipes and every pump --- exhaustively under the adopted
configuration (\texttt{--contingency-exhaustive-trunk}) rather than sampled ---
plus a bounded uniform sample of the remaining links, as singles and optional
trunk pairs, with non-converged solves discarded. Sources: \textbf{reservoirs
unbounded}, \textbf{tanks} at nominal delivered outflow (pump-fed tanks excepted,
Sec.~\ref{app:importer}). Functionality scale $N=3$; flow quantities rescaled by
$10^6$. \textbf{Priorities: none} --- every junction at the default tier, so the
allocation is a flat max-min fair share.

\emph{Engine configuration.} Tiered max-min fair-share allocation (the default),
integer fixed-point scale $10^3$, monotone fixed-point iteration to convergence.

\emph{Ground truth.} One steady-state pressure-driven solve per situation
(EPANET via WNTR, demand model PDD, required pressure $20\,$m, minimum pressure
$0\,$m, duration 0, the model's global demand pattern cleared, every pump opened
to match the working-condition baseline, broken links closed). Two corrections:
solves reporting ``system unbalanced'' are discarded, and demand junctions with
no undirected path to any source through in-service links are forced to ratio $0$
before quantization (those are nodes that cannot phisically receive water).

\emph{Situations.} Four families, each capped at $\sim\!60$ per network (10
single-, 20 double-, 30 triple-element closures where enough elements exist).
\emph{cluster}: breakable links within 2 topological hops of a random epicenter,
over seeds $\{1,2,3\}$ --- the one random family. \emph{targeted}: top-20\%-diameter
pipes plus all pumps, closed in singles/pairs/triplets, deterministic, largest
first. \emph{source}: reservoir-outlet closures. \emph{tank}: tank isolation
(every incident link). The combinatorial families are deterministic, so a network
with few reservoirs or tanks contributes proportionally fewer of those. Failed
ground-truth solves are skipped, leaving $n = 939$ scored situations.

\emph{Comparison semantics --- what makes it ``the same situation''.}
(i)~\textbf{Interventions are binary.} A situation closes links \emph{fully};
intermediate functionality levels never appear as benchmark \emph{inputs},
because EPANET has no partially degraded pipe. The module's proportional capacity
scaling is therefore exercised only in the \emph{outputs}, not validated as an
input mechanism. (ii)~\textbf{Node outages are incident-link closures}, since
EPANET cannot remove a node; on the module side the node itself is set to the
lowest level. (iii)~\textbf{Two different graphs are compared.} The module
propagates over the \emph{imported} model (inline pump/valve nodes, split
bidirectional edges --- not isomorphic to the \texttt{.inp} graph); junctions
match by round-trip identity, and each broken link maps to its imported edge(s)
and inline node. (iv)~\textbf{Scored population}: only demand-bearing junctions,
weighted by demand. (v)~\textbf{Out of scope}: EPANET control rules remain in the
ground-truth model (at a single steady state only their $t{=}0$ effect matters,
and we accept the result as ground truth); emitters, leakage, valve settings
beyond a capacity cap and minor losses are not represented; tank levels are fixed
at initial values, no drawdown. (vi)~\textbf{Thresholds}: levels quantize as
$\max(1,\min(N,\lceil r \cdot N\rceil))$, so at $N{=}3$ ``critical'' means served
ratio $\le 1/3$.

\emph{Metric aggregation.} Within a situation, FMS is a demand-weighted mean over
that situation's demand-bearing junctions; the reported figure is the mean of those
per-situation values. The $95\%$ intervals in the main text are percentile bootstrap
intervals over situations ($10^4$ resamples of the paired per-situation differences).

Detection is aggregated the same way: precision and recall are computed inside each
situation and then averaged. This makes each situation one unit of evidence --- one
failure, called correctly or not --- which is the question a resilience assessment
actually poses, and it gives a break on a nine-junction network the same weight as one
on a six-hundred-junction network.

Two situation classes need a stated convention. In $187$ of the $939$ situations the
network's redundancy absorbs the disruption entirely and no junction is critical in the
ground truth; recall is undefined there, and we credit a predictor that correctly
reports none with recall $1$. Symmetrically, a predictor that flags nothing anywhere
has undefined precision, and we credit it with precision $1$. The alternative is to
discard those situations, which would drop a fifth of the benchmark --- and precisely
the fifth where a false alarm is most costly, since the correct answer is that service
holds. An absorbed disruption is a question with an answer, and answering it correctly
is a success.

% ============================================================================
\section{Input Requirements by Mechanism}\label{app:inputs}

Table~\ref{tab:inputs} contrasts what a hydraulic simulation needs with what each
propagation mechanism needs. It is the basis of the architecture's data-parsimony
claim and of the cost argument in Sec.~\ref{app:cost}: the logical propagation needs
only structure and category assignments --- no demand, capacity, or hydraulic data
--- which is the structural reason it requires no physical oracle to validate
(main paper, Sec.~6), unlike the flow module.

\begin{table}[h]
\centering
\caption{Input requirements: hydraulic simulation vs.\ the two propagation
mechanisms. A hydraulic solver takes diameter, roughness, length, elevation and pump
curves and derives what a conduit carries; the flow module takes that carrying capacity
directly. Guards default to worst-case full dependency when undeclared, and shedding
priority is left unset in every experiment reported here}\label{tab:inputs}
\small
\begin{tabular}{@{}lccc@{}}
\toprule
Required input & EPANET/WNTR & Flow module & Logical propagation\\
\midrule
Topology (nodes, links) & \(\bullet\) & \(\bullet\) & \(\bullet\)\\
Category assignments & --- & \(\bullet\) & \(\bullet\)\\
Per-node demand & \(\bullet\) & \(\bullet\) & ---\\
Per-edge transport capacity & --- & \(\bullet\) & ---\\
Hydraulic specifications & \(\bullet\) & --- & ---\\
Shedding priority & --- & optional & ---\\
Guards (dependency level, backup) & --- & optional & optional\\
\bottomrule
\end{tabular}
\end{table}

% ============================================================================
\section{Extended Validation Across Failure Families}\label{app:extended}

Table~\ref{tab:extended} gives mean FMS for every network $\times$ family cell.
\emph{cluster} is the strongest family in aggregate ($0.972$) and \emph{source}
the weakest ($0.780$), with \emph{targeted} and \emph{tank} between them ($0.871$
each) --- but those two carry different content: \emph{targeted} is the
capacity-loaded family that forces the most rerouting, while \emph{tank} is
dominated by the benchmark-semantics artifact of Sec.~\ref{app:negative}.

The weakest well-sampled cell is Modena under \emph{targeted} ($0.636$ over $60$
situations); Net3's \emph{source} cell is lower still ($0.372$) but rests on three.
Absolute FMS is the wrong lens for comparing networks, though, since a network that barely degrades
under our families leaves almost no error to remove: summed shortfall
($1-\mathrm{FMS}$) tracks how many situations a network yields (Spearman $+0.76$
against situation count) more than how well the module does on it. Table~\ref{tab:skill}
therefore reports \emph{skill} against the null model.

\begin{table}[h]
\centering
\caption{Per-network relative performance. $\bar{e}$ is the mean per-situation
shortfall $1-\mathrm{FMS}$; $\bar{e}_{\mathrm{null}}$ the same for the all-operational
null model; skill is $1-\bar{e}/\bar{e}_{\mathrm{null}}$ ($1$ perfect, $0$ no better
than predicting no loss, negative worse). Ordered by skill}\label{tab:skill}
\small
\begin{tabular}{@{}lrrrrr@{}}
\toprule
Network & junctions & situations & $\bar{e}$ & $\bar{e}_{\mathrm{null}}$ & skill\\
\midrule
Net2       &  35 & 113 & 0.000 & 0.247 & $1.00$\\
Net1       &   9 &  68 & 0.003 & 0.439 & $0.99$\\
Aqueduct A & 417 & 123 & 0.007 & 0.310 & $0.98$\\
Aqueduct B & 439 &  75 & 0.018 & 0.049 & $0.64$\\
Aqueduct C & 607 & 118 & 0.132 & 0.216 & $0.39$\\
Modena     & 268 & 134 & 0.200 & 0.268 & $0.25$\\
C-Town     & 388 & 178 & 0.089 & 0.020 & $-3.40$\\
Net3       &  92 & 130 & 0.125 & 0.024 & $-4.21$\\
\bottomrule
\end{tabular}
\end{table}

The ordering this gives differs from the absolute one and is the more useful of the
two. Modena and Aqueduct~C carry the largest absolute shortfall yet still remove a
quarter to two-fifths of the null's error. Net3 and C-Town, middling in absolute
terms, are the two networks where the module is \emph{worse than doing nothing}: both
have $\bar{e}_{\mathrm{null}}$ near $0.02$, so their families barely disturb service at
all, and the module's over-warning costs several times the little there was to gain.
The ratio is volatile at so small a denominator --- the sign is not, and the sign is
the finding. The causes differ by network and are separated in the main paper's Sec.~6;
per-network critical-class precision and recall appear in Table~\ref{tab:capmethod}'s
``Uniform'' columns, that being the adopted configuration.

\begin{table}[t]
\centering
\caption{Mean FMS by network (rows) and failure family (columns), pooled over
seeds. The three bottom rows score the main paper's baselines on the identical
situations. A cell is empty when the network lacks the relevant element. Same
$n=939$ scored situations as the main paper's Table~1}\label{tab:extended}
{%
\begin{tabular}{@{}lcccc@{}}
\toprule
Network & cluster & targeted & source & tank\\
\midrule
Net1        & 1.000 & 0.987 & 1.000 & 0.909\\
Net2        & 1.000 & 1.000 & ---   & 1.000\\
Net3        & 0.917 & 0.842 & 0.372 & 1.000\\
Aqueduct A  & 0.995 & 0.991 & 1.000 & ---  \\
Aqueduct B  & 0.999 & 0.959 & 1.000 & 0.929\\
Aqueduct C  & 0.968 & 0.762 & 0.754 & 1.000\\
Modena      & 0.970 & 0.636 & 0.777 & ---  \\
C-Town      & 0.938 & 0.944 & 0.960 & 0.845\\
\midrule
Aggregate     & 0.972 & 0.871 & 0.780 & 0.871\\
Null baseline & 0.875 & 0.744 & 0.646 & 0.953\\
Reachability  & 0.969 & 0.855 & 0.729 & 0.977\\
\bottomrule
\end{tabular}%
}
\end{table}

The binary metric inverts the FMS parity of the bottom rows. Reachability's
precision is $1.00$ by construction --- it only ever calls a fully severed
junction critical, and severed junctions are critical by the shared ground-truth
correction --- so its recall \emph{is} its content, and that recall collapses
exactly where capacity matters: $0.54$ under \emph{targeted} attack against the
module's $0.90$, and $0.54$ under \emph{source} failure against $0.92$. Every
critical junction it misses is connected to a source yet starved. The module pays
for that recall in precision ($0.73$ on average), allocating scarce supply by fair
share and so over-predicting. All figures use the adopted default: a flat
allocation, priority unset throughout.

% ============================================================================
\section{What Carries the Fidelity: Orientation, Capacity, Priority}\label{app:ablation}

Three inputs shape the imported flow module: which way each edge points, how much each
pipe can carry, and who is protected under shortage. This section measures what each
contributes. They are different kinds of question. Orientation is an \emph{ablation}:
it is part of the shipped importer, so it can be removed and the loss measured.
Capacity is a \emph{method comparison}: two candidate ways of setting it, one of which
ships. Priority is a \emph{capability measurement}: what the input can make the model
do, a separate question from whether it improves agreement with hydraulics, since a
protect-ordering has no counterpart in WNTR to be checked against. All run under the
protocol of Sec.~\ref{app:repro} with the canonical configuration otherwise unchanged
--- orientation and capacity over all eight networks, priority over two.

\emph{Orientation is what the import-time solves buy.} Under the shipped configuration
the demand and contingency solves serve one purpose: they \emph{orient} edges from the
sign of the flow, including pipes that carry flow only once a closure reroutes onto
them. Capacity comes from a constant and priority is left unset, so nothing else in the
pipeline consumes them. Table~\ref{tab:orient} removes orientation and measures the
cost: in the ablated arm every pipe becomes two opposed directed edges, the graph as it
stands with no direction information at all. Capacity is identical between arms, the
uniform design velocity never reading the solves.

Orientation buys \emph{recall} and costs \emph{precision}. Discarding it moves average
recall from $0.944$ to $0.898$ and average precision from $0.729$ to $0.862$: the extra
edges let flow reroute along paths the real network does not offer, so fewer junctions
are called starved --- which removes false alarms and conceals real ones together.
Recall falls on five of the eight networks, by as much as $0.239$. A
resilience assessment is judged on the criticals it finds, and that is why the importer
orients.

One network settles the question by itself. Aqueduct~C is the only one where removing
orientation costs on every measure at once --- FMS $-0.144$, precision $-0.136$, recall
$-0.239$ --- and it is the most strongly directional network we import. There the
phantom reroutes do not merely soften the module's answer; they invent supply routes
that carry the network through failures it does not survive.

The same table localizes C-Town's over-warning. Its precision rises from $0.130$ to
$0.774$ once orientation is discarded, the largest single effect here, which places the
cause in the orientation the sweep assigns rather than in capacity, supply or
allocation. C-Town is the network whose seven tanks are all pump-fed, so its flow
directions genuinely reverse as pumps cycle, and one orientation drawn from a single
operating regime is the wrong object to give it. Why the sweep settles on the
directions it does there is open; that the over-warning follows from them is not.

\begin{table}[t]
\centering
\caption{Orientation ablation over all eight networks. ``Oriented'' is the shipped
importer, taking each pipe's direction from the sign of its simulated flow;
``Undirected'' splits every pipe into two opposed directed edges, the graph with no
direction information. Capacity, demand, priority, situations and ground truth are
identical between arms. \emph{edges} is the oriented edge count. Bold marks the better
arm where they differ}\label{tab:orient}
\small
\begin{tabular}{lrccc@{\hskip 1.5em}ccc}
\toprule
 & & \multicolumn{3}{c}{Oriented (shipped)} & \multicolumn{3}{c}{Undirected} \\
\cmidrule(lr){3-5}\cmidrule(l){6-8}
Network & edges & FMS & $P$ & $R$ & FMS & $P$ & $R$ \\
\midrule
Net1       &   20 & 0.997 & 1.000 & 1.000 & 0.997 & 1.000 & 1.000 \\
Net2       &   60 & 1.000 & 1.000 & 1.000 & 1.000 & 1.000 & 1.000 \\
Net3       &  194 & 0.875 & 0.775 & \textbf{0.991} & \textbf{0.882} & \textbf{0.849} & 0.977 \\
Aqueduct A &  771 & 0.993 & 0.979 & \textbf{0.942} & \textbf{0.994} & \textbf{1.000} & 0.937 \\
Aqueduct B &  853 & 0.982 & 0.949 & \textbf{0.855} & \textbf{0.987} & \textbf{0.993} & 0.723 \\
Aqueduct C & 1290 & \textbf{0.868} & \textbf{0.874} & \textbf{0.890} & 0.724 & 0.738 & 0.651 \\
Modena     &  575 & \textbf{0.800} & 0.635 & 0.972 & 0.787 & \textbf{0.714} & \textbf{0.974} \\
C-Town     &  655 & \textbf{0.911} & 0.130 & \textbf{0.908} & 0.869 & \textbf{0.774} & 0.889 \\
\midrule
Average    &      & \textbf{0.918} & 0.729 & \textbf{0.944} & 0.892 & \textbf{0.862} & 0.898 \\
\bottomrule
\end{tabular}
\end{table}

\emph{Capacity method.} The importer sizes every pipe at a single design
velocity, $\tfrac{\pi}{4}d^2 v_{\text{design}}$ with $v_{\text{design}}=2.5$\,m/s. The
alternative is a per-pipe ``drill'' that sizes each pipe from the peak velocity it
reaches across the hydraulic solves, on the reasoning that a busier pipe should carry
more. Table~\ref{tab:capmethod} runs both arms with orientation held fixed.

The constant wins. It gives higher FMS on five networks and matches the drill on the
other three, higher precision on four and lower on one, and averages $0.918$ FMS against
$0.905$ and $0.729$ precision against $0.707$, for $F_1$ $0.823$ against $0.813$. The gains are largest on
the capacity-loaded networks (precision $+0.08$ on Net3, $+0.07$ on Aqueduct~C). The
drill buys back a little recall ($0.957$ against $0.944$ on average, from flagging more
junctions) and C-Town's already-degenerate precision is marginally higher under it
($0.133$ against $0.130$). We ship the constant: it matches or beats a hydraulic
discovery step while being simpler, standard, and reproducible without any solve. It
must be sized correctly, though --- at $1$--$2$\,m/s a constant is too small and scores
worse (Sec.~\ref{app:negative}).

\begin{table}[t]
\centering
\caption{Capacity-method ablation, four families pooled per network. ``Uniform'' =
capacity from a single $2.5$\,m/s design velocity (adopted); ``Drill'' brings per-pipe
peak-velocity discovery from the hydraulic solves. Orientation is identical in both
arms and both leave priority unset. $P$/$R$ are critical-class precision and recall.
Bold marks the better arm where they differ}\label{tab:capmethod}
\small
\begin{tabular}{lccc@{\hskip 1.5em}ccc}
\toprule
 & \multicolumn{3}{c}{Uniform (adopted)} & \multicolumn{3}{c}{Drill} \\
\cmidrule(r){2-4}\cmidrule(l){5-7}
Network & FMS & $P$ & $R$ & FMS & $P$ & $R$ \\
\midrule
Net1       & 0.997 & 1.000 & 1.000 & 0.997 & 1.000 & 1.000 \\
Net2       & 1.000 & 1.000 & 1.000 & 1.000 & 1.000 & 1.000 \\
Net3       & \textbf{0.875} & \textbf{0.775} & 0.991 & 0.843 & 0.700 & \textbf{0.997} \\
Aqueduct A & 0.993 & 0.979 & 0.942 & 0.993 & 0.979 & 0.942 \\
Aqueduct B & \textbf{0.982} & \textbf{0.949} & 0.855 & 0.967 & 0.940 & 0.855 \\
Aqueduct C & \textbf{0.868} & \textbf{0.874} & 0.890 & 0.845 & 0.803 & \textbf{0.928} \\
Modena     & \textbf{0.800} & \textbf{0.635} & 0.972 & 0.778 & 0.616 & \textbf{0.989} \\
C-Town     & \textbf{0.911} & 0.130 & 0.908 & 0.902 & \textbf{0.133} & \textbf{0.934} \\
\midrule
Average    & \textbf{0.918} & \textbf{0.729} & 0.944 & 0.905 & 0.707 & \textbf{0.957} \\
\bottomrule
\end{tabular}
\end{table}

\emph{Priority as controllability.} Shedding priority states whom to protect. It is a
normative input, and WNTR has no equivalent of it, so we measure whether the model
\emph{realizes} a specified protect-ordering rather than whether that ordering agrees
with hydraulics. For each situation we take a target shed/protect pattern --- that
situation's own WNTR-served set serves as a well-defined one --- and ask two questions.
How far can a priority vector steer the model toward the target? And does one fixed
vector keep that effect on situations it was not built from? Table~\ref{tab:steering}
answers both on Modena and C-Town, chosen because between them they show the two kinds
of divergence the module produces. The mechanism throughout is the allocation's
strict-tier preemption: a higher-priority junction is rationed before any
lower-priority one competes.

\begin{table}[h]
\centering
\caption{Shedding-priority steering: demand-weighted agreement with the target pattern,
before $\to$ after, with the gain in bold where material. \emph{Per-situation} refits a
vector for every situation and measures the reachable range. \emph{Single vector} fits
one vector per network on its ten worst situations, then scores it both there and on
the remaining situations it was not fitted to. Both vectors are built from the ground
truth, so neither is available to the importer (Sec.~\ref{app:negative})}\label{tab:steering}
\small
\begin{tabular}{@{}lcc@{}}
\toprule
 & Modena & C-Town\\
\midrule
per-situation, worst-10     & $0.466 \to 0.831$ \textbf{(+0.365)} & $0.590 \to 0.679$ (+0.089)\\
\quad of which moved        & \textbf{10 of 10} & 1 of 10\\
single vector, fit set      & $0.466 \to 0.788$ \textbf{(+0.322)} & $0.590 \to 0.657$ (+0.067)\\
single vector, held out     & $0.827 \to 0.907$ \textbf{(+0.080)} & $0.930 \to 0.930$ (+0.000)\\
\bottomrule
\end{tabular}
\end{table}

\emph{The range is wide, and bounded by the kind of error.} On Modena all ten of the
worst situations improve, each by $+0.33$ to $+0.55$. On C-Town one situation improves
enormously (\texttt{cluster\#36}, $0.040 \to 0.904$) and the remaining nine, all from
the \emph{tank} family, move by $+0.012$ or less --- and that is under a per-situation
oracle, the most favourable condition the lever can be given. The reason is that
priority \emph{redirects} a shortage; it cannot correct one. Modena's divergence is
distributional: fair share spreads a shortage that pressure-driven hydraulics
concentrate, so restating who is protected repairs it. C-Town's tank divergence is the
pump-fed over-warning of Sec.~\ref{app:negative}, where the shortage itself is wrong,
and no ordering of the victims helps. Where the total is right and the distribution
wrong the lever has wide range; where the total is wrong it has none.

\emph{Generalization splits the same way.} One fixed vector per network keeps most of
the range on Modena, including $+0.08$ on situations it was never fitted to, and none
of it on C-Town. So on a network whose errors are distributional, a single expert-set
ordering transfers beyond the situations used to choose it.

This does not make priority derivable. Both vectors are fitted from the ground truth
the importer is trying to predict, and the orderings that can actually be derived from
a network --- solve-free proxies and a contingency-severity ranking --- fail
(Sec.~\ref{app:negative}). The claim is an authoring one: an expert who knows how a
utility distributes water under shortage states it in one attribute per node, and on a
distributional network that statement both realizes the intended policy and carries to
situations it was not tuned on. Both arms:
\texttt{experiments/aqueducts/priority\_steering.py}.

% ============================================================================
\section{Ground-Truth Threshold Sensitivity}\label{app:sensitivity}

The comparison inherits exactly one hydraulic constant --- the PDD
service-pressure convention (required $20$\,m, minimum $0$\,m) --- and one
quantization choice, the scale size $N{=}3$. A headline that held only at one
arbitrary setting of either would be weaker than it looks, so we re-ran the full
protocol on Net1 and Net3 (seeds $\{1,2,3\}$, $192$ scored situations per arm)
varying required pressure $\{15,20,25\}$\,m, minimum pressure $\{0,5\}$\,m, and
$N \in \{3,5,10\}$.

Mean FMS spans $0.988$--$0.992$ across every arm against that sweep's own
baseline of $0.989$ --- a spread of $0.004$, far below the effect of any
modelling choice tested elsewhere. The $N{=}10$ arm, where agreement must hold at
a granularity three times finer than the headline scale, loses only $0.001$: the
module's delivered-ratio estimates track the hydraulic ratios themselves, not
just their coarse bucketing. Neither the service-pressure convention nor the
quantization granularity is doing the work. (This sweep predates the
uniform-capacity configuration, so its absolute level sits above the main paper's
Table~1; what it establishes is the \emph{insensitivity}, a within-sweep
comparison.)

% ============================================================================
\section{Negative Results: What Did Not Improve Fidelity}\label{app:negative}
% Source: experiments/aqueducts/ATTEMPTS.md (2026-07-12..14) and
% archive/pre-rebuild-variants/{algorithm_variants,importer_variants}.py.
% Benchmark: the 10 lowest-FMS situations across five networks (worst-10),
% ground truth corrected for the singular-PDD artifact described below.
We record the alternatives we rejected, since several are the ``obvious'' fixes a
reader might propose. Unless stated otherwise these were measured on the ten situations
where the module diverges most from the ground truth: mean FMS $0.770$ for the pipeline
as first described, which the capacity-margin and pipe-splitting refinements raise to
$0.927$. These are historical measurements taken on that earlier pipeline during the
exploration logged in \texttt{experiments/aqueducts/ATTEMPTS.md}, which records each
figure and the run it came from; the shipped configuration is the one every other
section reports, and re-running these arms today would need the superseded importer.

\emph{Capacity-margin sensitivity (retained drill only).} The $\times 2$ margin
under the per-pipe drill --- the adopted default uses a fixed design velocity and
no margin --- is an exposed parameter, not a value fitted to the evaluation.
Swept over $\{1, 1.25, 1.5, 1.75, 2, 3\}$, fidelity is \emph{insensitive} on four
networks (Net1, Net2, Aqueducts~A and~B, flat to two decimals) and climbs to a
plateau on two (Net3 $0.95\!\to\!1.00$, knee near $\times 1.5$; Aqueduct~C
$0.14\!\to\!0.89$ by $\times 2$). Pooled FMS rises $0.83\!\to\!0.96$ from
$\times 1$ to $\times 2$ and only $+0.01$ more by $\times 3$; recall is flat
throughout while precision tracks FMS. So $\times 2$ sits on the plateau for
every network, and the velocity it implies stays under the $3$\,m/s design ceiling the
importer caps every margined pipe at. Held-out selection on other
public networks (ky10, ky4, Net6) was attempted and abandoned --- they do not
converge or self-starve at peak demand without per-network calibration, and
skeletonizing does not help. Harness:
\texttt{experiments/aqueducts/margin\_sweep.py}.

\emph{Source supply.} Setting \emph{reservoir} supply to a bounded nominal outflow rate
(as we do for tanks) breaks source-failure validation: a failed reservoir's
survivors, now capped, cannot cover its lost demand, so the module over-predicts
wholesale, whereas EPANET's fixed-head reservoirs supply the slack. Reservoirs
are therefore left unbounded.

The symmetric effect appears on \emph{gravity} tanks, which we do cap at
sustainable nominal outflow while the $t=0$ ground truth treats them as
fixed-head nodes draining as fast as their pipes permit. On the Net3 situation
where the module diverges most (\texttt{cluster\#25}, FMS $0.058$, which no
importer or allocation lever moves) a break isolates a district onto three local
gravity tanks. A four-way max-flow discriminator localizes the binding constraint
to source supply rather than conveyance or orientation: directed with real supply
$48.5\%$ of demand served, directed with supply $=\infty$ $84.9\%$, bidirectional
with real supply $49.9\%$, bidirectional with supply $=\infty$ $100\%$. Measured
against the cap, WNTR draws $6.6\times$ the nominal rate from those tanks (one at
$20\times$), giving $48\%$ served against the WNTR solve's $134\%$. This is a
depletion-aware model read against a depletion-blind snapshot; an extended-period
ground truth would narrow it. We keep the cap regardless: raising it to match the
$t=0$ surge would make the module optimistic about supply no real tank sustains.

This is \emph{not} the explanation for the tank family's average precision of
$0.20$ (main paper, Sec.~6). That number is C-Town, which supplies $57$ of the
family's $72$ situations and scores $0.00$ on every one; its tanks are pump-fed
pass-throughs sized from incident pipes, so no cap binds them. Their over-warning
follows instead from the orientation the sweep assigns: discarding orientation lifts
C-Town's precision from $0.130$ to $0.774$ (Table~\ref{tab:orient}). We separate the
two mechanisms rather than let one stand in for the other.

\emph{Deriving a shedding order.} Three solve-free proxies for junction
vulnerability (shortest Hazen-Williams distance to a source, junction-to-source
effective resistance, and effective resistance under topological contingency)
each scored \emph{worse} than a flat allocation. A more elaborate
contingency-solve ranking --- closing cycle-trunk links singly, in pairs and in
triplets, then ordering junctions by accumulated demand-weighted unmet service
--- shifts the operating point but only as a trade (Table~\ref{tab:priority}):
recall $0.90 \to 0.92$, precision $0.70 \to 0.67$, FMS one point. Which junctions
a network abandons under structural failure is a demand-and-pressure property no
static ordering captures, and an ordering is in any case operator policy rather
than prediction. The importer therefore assigns none.

\begin{table}[h]
\centering
\caption{Why the importer derives no shedding order: critical-class detection
over the targeted and source families, flat allocation versus the
cycle-aware contingency severity ranking. The ranking trades precision for recall
rather than improving agreement}\label{tab:priority}
\begin{tabular}{@{}lccc@{}}
\toprule
Shedding order & precision & recall & FMS\\
\midrule
none (flat, as reported)     & 0.70 & 0.90 & 0.87\\
contingency severity ranking & 0.67 & 0.92 & 0.88\\
\bottomrule
\end{tabular}
\end{table}

\emph{Allocation shape.} The control arm was the single-solve min-cost max-flow
(the module's alternative strategy). Replacing it with global proportional
water-filling --- every consumer served the same feasible fraction --- is sharply
worse ($0.373$): uniform allocation punishes consumers unrelated to the local
bottleneck, which real hydraulics never does. Max-min fair share does beat the LP
on \emph{as-imported} capacities ($0.774$ vs $0.724$, never losing), but the
entire advantage vanishes once the capacity margin is applied ($0.918$ vs
$0.927$), at materially more solver work. The apparent allocation-shape problem
was an importer-attribute problem --- which is why the paper treats fair share as
a scarcity-sharing \emph{policy} choice (equal-priority consumers degrade
together instead of one being arbitrarily zeroed) rather than a fidelity
optimization.

\emph{Orientation and component fixes.} Adding a reverse edge to every one-way
pipe (making the flow graph effectively undirected, as EPANET's links are)
changes essentially nothing ($0.772$) and occasionally hurts via degenerate ties:
the sweep's confident one-way calls are reliable, and the real directional defect
was only the \emph{proportional} capacity split on pipes already known to be
bidirectional. Lifting pump-curve and valve capacity caps, and removing source
supply limits entirely, each had zero effect on all ten situations --- neither is
ever the binding constraint. Replacing the sweep velocities with a low constant
($1$--$2$\,m/s) looked worse ($0.681$--$0.791$), which we first read as evidence
that the sweep's relative capacity profile carries information; the full ablation
of Sec.~\ref{app:ablation} overturns that, the earlier result being an artifact
of an under-sized constant on two networks.

\emph{A dropped family.} \emph{Partial, system-wide} source degradation was
tried and dropped: when every source weakens together the loss is diffuse and a
trivial ``flag every junction'' baseline matches the module. It was replaced by
concentrated source \emph{failure} (outlet closure).

\emph{A ground-truth artifact, not a module error.} Where breaks sever a
component from every source, the PDD system for that component is singular and
WNTR converges, without warning, to an arbitrary internal circulation whose
per-junction delivered demands read as fully served while summing to $\approx 0$.
Two of the ten worst situations owed their entire apparent divergence to this:
the module's ``everyone critical'' answer was correct and the ``ground truth''
was not. We post-process the WNTR solve, marking demand junctions with no
undirected path to a source through in-service elements as unserved, and report
all FMS against the corrected reference.

% ============================================================================
\section{Computational Cost}\label{app:cost}

Measured on the same single machine as every other number here (consumer Linux
workstation, single core, Python 3.12; absolute times are machine-shaped, ratios
are the content). Harness: \texttt{experiments/aqueducts/cost\_benchmark.py}.

\emph{One-time import.} Parse plus the hydraulic sweeps that orient edges cost
$0.1$--$5.9$\,s per network under the adopted exhaustive-trunk configuration.
This is paid once per network, not per situation, and it is the only place the
importer solves hydraulics at all.

\emph{Per-situation.} Table~\ref{tab:cost} times one engine Propagation against
one WNTR PDD solve of the same situation. The engine figure includes the deep copy of
the model the harness takes before mutating it; we time that copy separately but leave
it in the reported number rather than subtract it. Which situations are timed decides the
answer, because engine cost tracks how much \emph{scarcity} a situation creates
rather than how large the network is: the fair-share water-filling bisects for a
common allotment and pays repeated max-flow solves per round, so a break that
merely disconnects a leaf costs almost nothing while one forcing a large
contended reroute costs seconds. We therefore sample evenly across the four
families rather than in proportion to their sizes, which would let the mild
\emph{cluster} family (459 of the 939 situations) set the figure. The
within-network spread this exposes is part of the result: on Net3 the slowest
situation is nearly $100\times$ the median.

\begin{table}[h]
\centering
\caption{Per-situation cost, median over $3$ repeats of situations drawn evenly
from the four families, observed range in brackets. Ratio is of the medians. The
engine is at parity on textbook networks and one to two orders of magnitude
slower at aqueduct scale}\label{tab:cost}
\small
\begin{tabular}{@{}lrrrr@{}}
\toprule
Network & import & engine (ms) & WNTR (ms) & ratio\\
\midrule
Net1       & $0.1$\,s & $12.6$ \; [$1.9$--$100$]        & $8.4$   & $1.5\times$\\
Net3       & $0.4$\,s & $51.4$ \; [$41$--$4907$]        & $23.2$  & $2.2\times$\\
Aqueduct A & $2.8$\,s & $4674$ \; [$3880$--$9674$]      & $82.2$  & $57\times$\\
Aqueduct B & $4.3$\,s & $11152$ \; [$589$--$13687$]     & $104.4$ & $107\times$\\
Aqueduct C & $5.9$\,s & $17039$ \; [$4083$--$25283$]    & $129.3$ & $132\times$\\
\bottomrule
\end{tabular}
\end{table}

The conclusion is not a favourable one, and we state it plainly: this paper
claims no runtime advantage over a dedicated hydraulic simulator, and at aqueduct
scale the gap is not a constant factor to be tuned away but one to two orders of
magnitude, with individual situations taking tens of seconds. 
The architecture's case is the input requirements of
Sec.~\ref{app:inputs} --- what it models without hydraulic data, and what it
composes --- not speed. 

\end{document}
